\documentclass[12pt,a4paper]{report}
\usepackage[utf8]{inputenc}
\usepackage[T1]{fontenc}
\usepackage[margin=2.5cm]{geometry}
\usepackage{fancyhdr}
\usepackage{hyperref}

\pagestyle{fancy}
\fancyhf{}
\fancyhead[L]{\leftmark}
\fancyhead[R]{\thepage}

\begin{document}
\begin{titlepage}
    \centering
    {\Large\textbf{PROGRAMMING REPORT}}\\[2cm]
    {\Huge\textbf{Hello World Example}}\\[3cm]
    \textbf{Prepared by:} Programmer Name\\[0.5cm]
    {\large Date: \today}
\end{titlepage}

\tableofcontents
\newpage

\chapter{Introduction}
The purpose of this report is to demonstrate a simple programming concept using the printf function to display "hello world" on the screen.

\chapter{Methodology}
To achieve this, we will use a programming language such as C, which supports the printf function. The code will be as follows: 
\begin{verbatim}
#include <stdio.h>
int main() {
    printf("hello world");
    return 0;
}
\end{verbatim}

\chapter{Results}
When the above code is compiled and executed, the output will be "hello world" printed on the screen.

\chapter{Conclusion}
In conclusion, the printf function is a basic yet essential part of programming in C, allowing users to print output to the screen. This example demonstrates its simplicity and effectiveness.

\end{document}